\chapter*{F}
\label{chap:f}

\term{Faà di Bruno's Formula}
The general chain rule for higher derivatives: $\frac{d^n}{dx^n}f(g(x)) = \sum \frac{n!}{k_1!(m_1!) \cdots k_n!(m_n!)} f^{(m_1+\cdots+m_n)}(g(x)) \prod_{j=1}^n \left(\frac{g^{(j)}(x)}{j!}\right)^{m_j}$ where the sum is over all non-negative integers $k_1, \ldots, k_n, m_1, \ldots, m_n$ satisfying $\sum j k_j = n$ and $\sum k_j = m$. This formula appears in combinatorics, probability (Bell polynomials), and differential equations.

\term{Face (Polytope)}
A face of a convex polytope is the part cut out by a supporting hyperplane. Faces include vertices, edges, facets, and the whole polytope. \par\noindent\textbf{Applications:} faces describe the geometry of polyhedra and are used in linear programming, computer graphics, and computational geometry.

\term{Factorial Function}
The factorial is $n!=1\cdot2\cdots n$ for a nonnegative integer, with $0!=1$. Stirling's approximation gives $n!\sim\sqrt{2\pi n}(n/e)^n$, and $n!=\Gamma(n+1)$. \par\noindent\textbf{Applications:} factorials count arrangements, enter probability formulas, and support series expansions and numerical approximations.

\term{Fadeev--Popov Determinant}
In gauge theory, the Fadeev--Popov determinant arises when fixing a gauge in the path integral. For a gauge group $G$ acting on a field configuration space, the Fadeev--Popov procedure inserts $1 = \Delta_{FP}[A] \int \mathcal{D}g \, \delta(G(A^g))$ into the path integral. The resulting ghost fields are anticommuting scalars and are essential for maintaining unitarity in non-abelian gauge theories.

\term{Fan (Toric Geometry)}
A fan $\Sigma$ in $\mathbb{R}^n$ is a collection of strongly convex rational polyhedral cones such that every face of a cone in $\Sigma$ is also in $\Sigma$, and the intersection of any two cones is a face of each. Fans correspond to toric varieties: each fan $\Sigma$ defines a complex algebraic variety $X_\Sigma$ with an effective torus action. The fan for $\mathbb{CP}^n$ consists of cones generated by subsets of the standard basis and their negatives.

\term{Fan Graph}
A fan graph joins one central vertex to every vertex of a path. \par\noindent\textbf{Applications:} fan graphs model hub-and-spoke networks and appear in chemical graph theory, connectivity studies, and graph spectra.

\term{Fatou's Lemma}
For a sequence of non-negative measurable functions $f_n \geq 0$: $\int \liminf f_n \, d\mu \leq \liminf \int f_n \, d\mu$. This fundamental result in measure theory is a key ingredient in the proof of the Dominated Convergence Theorem and the Monotone Convergence Theorem. It reflects the lower semicontinuity of the integral under pointwise convergence.

\term{Fault-Tolerant Quantum Computation}
A fault-tolerant quantum computation uses quantum error-correcting codes to protect against errors during computation. The threshold theorem states that if the error rate per gate is below a constant threshold, arbitrary long quantum computations can be performed with polynomial overhead. Surface codes and the $[[7,1,3]]$ Steane code are prominent examples. Fault tolerance requires entanglement, syndrome measurement, and adaptive correction.

\term{Fermat's Last Theorem}
Fermat's Last Theorem states that $x^n+y^n=z^n$ has no positive integer solutions when $n>2$. \par\noindent\textbf{Applications:} its proof connects elliptic curves, modular forms, and number theory and illustrates how elementary equations can require deep mathematics.

\term{Fermat's Little Theorem}
If $p$ is prime and $p\nmid a$, then $a^{p-1}\equiv1\pmod p$. \par\noindent\textbf{Applications:} Fermat's little theorem supports primality tests, modular inverses, and public-key cryptography.

\term{Fiber Bundle}
A fiber bundle is a space $E$ with a projection $\pi:E\to B$ that locally looks like a product $U\times F$, where $B$ is the base and $F$ is the fiber. Globally, the fibers may twist, as in a Möbius strip. \par\noindent\textbf{Applications:} bundles model fields, orientations, configuration spaces, and gauge potentials in geometry and physics.

\term{Fibrant Object}
In a model category, an object is fibrant when its map to the terminal object has the required lifting property. Informally, a fibrant object is one on which homotopy-theoretic constructions behave smoothly. In simplicial sets, the fibrant objects are Kan complexes. Weak equivalences between objects that are both fibrant and cofibrant have homotopy inverses. \par\noindent\textbf{Applications:} fibrant replacements let algebraic models represent the homotopy type of spaces without changing their essential information.

\term{Fibration}
A fibration is a map with a homotopy lifting property: a deformation in the base can be lifted after an initial lift is chosen. Fiber bundles are common examples. \par\noindent\textbf{Applications:} fibrations organize homotopy groups, fibered spaces, and the deformation of geometric or topological models.

\term{Fibonacci Sequence}
The Fibonacci sequence satisfies $F_0=0$, $F_1=1$, and $F_n=F_{n-1}+F_{n-2}$. The ratio of consecutive terms approaches the golden ratio $\phi=(1+\sqrt5)/2$. \par\noindent\textbf{Applications:} Fibonacci numbers appear in combinatorics, algorithms, continued fractions, coding, and models of branching and growth.

\term{Field (Algebra)}
A field is a commutative ring in which every nonzero element has a multiplicative inverse. Examples include $\mathbb Q$, $\mathbb R$, $\mathbb C$, and finite fields. \par\noindent\textbf{Applications:} fields provide the scalars for linear algebra and are central to coding theory, cryptography, and polynomial equations.

\term{Field Extension}
A field extension $L/K$ is a pair of fields with $K \subseteq L$. The degree $[L:K]$ is the dimension of $L$ as a $K$-vector space. A finite extension is algebraic if every element of $L$ is a root of a polynomial over $K$. The splitting field of a polynomial $f$ over $K$ is the smallest extension in which $f$ factors into linear factors. Galois theory studies the correspondence between intermediate fields and subgroups of the Galois group.

\term{Fierz Identity}
In quantum field theory, the Fierz identity is a rearrangement formula for products of spinor bilinears. For Dirac spinors $\psi_i$: $(\bar{\psi}_1 \Gamma^A \psi_2)(\bar{\psi}_3 \Gamma_A \psi_4) = \sum_B c_B (\bar{\psi}_1 \Gamma^B \psi_4)(\bar{\psi}_3 \Gamma_B \psi_2)$ where the sum runs over the 16 Dirac matrices and $c_B$ are numerical coefficients. Fierz identities are used in deriving four-fermion interactions and checking supersymmetry transformations.

\term{Filtration}
A filtration is a nested sequence $\mathcal F_0\subseteq\mathcal F_1\subseteq\cdots$ that records how an object is built in stages. In probability, a filtration represents the information available over time; in algebra, it records layers of a module or complex. \par\noindent\textbf{Applications:} filtrations are used in stochastic processes, spectral sequences, persistent homology, and multiscale data analysis.

\term{Final Object}
In category theory, a final (or terminal) object is an object $T$ such that for every object $X$ there exists a unique morphism $X \to T$. In the category of sets, the final object is any singleton $\{*\}$. In the category of groups, it is the trivial group. Final objects are unique up to unique isomorphism. Dually, an initial object has a unique morphism from it to every object.

\term{Finite Field}
A finite field is a field with finitely many elements; its cardinality is necessarily a prime power $p^n$ and its characteristic is $p$. For every prime power $q = p^n$ there exists a field $\mathbb{F}_q$ with $q$ elements, unique up to isomorphism, namely the splitting field of $x^q - x$ over $\mathbb{F}_p$. The multiplicative group $\mathbb{F}_q^\times$ is cyclic, and a generator is a primitive element. The Frobenius automorphism $x \mapsto x^p$ generates the Galois group of $\mathbb{F}_q$ over $\mathbb{F}_p$. Finite fields underlie coding theory (Reed--Solomon and BCH codes), cryptography (elliptic curves, discrete logarithms), and pseudorandom generation.

\term{Finite Group}
A finite group has finitely many elements, and its order is the number of elements. \par\noindent\textbf{Applications:} finite groups describe symmetries of molecules, crystals, codes, puzzles, and finite computational systems.

\term{Finite Set}
A finite set has a fixed number of elements and is in bijection with $\{1,2,\ldots,n\}$ for some $n$. \par\noindent\textbf{Applications:} finite sets model stored data, finite state machines, discrete experiments, and combinatorial objects.

\term{Finite-difference Method}
A finite-difference method approximates derivatives by differences of function values on a grid. \par\noindent\textbf{Applications:} it solves heat, wave, fluid, and structural equations numerically.

\term{Finite-element Method}
A finite-element method approximates a PDE solution by piecewise polynomial functions on a mesh, usually through a variational formulation. \par\noindent\textbf{Applications:} it is widely used for stress, heat, fluid, electromagnetic, and structural simulations.

\term{Finite-volume Method}
A finite-volume method integrates a conservation law over small control volumes and balances fluxes across their boundaries. \par\noindent\textbf{Applications:} its conservation property makes it effective for computational fluid dynamics, weather, combustion, and traffic flow.

\term{Finitely Presented Group}
A finitely presented group is given by generators and relations: $G = \langle S \mid R \rangle$ where $S$ is a finite set of generators and $R$ is a finite set of relations (words in $S$). Every finite group is finitely presented. However, the word problem is undecidable in general: there is no algorithm to determine whether two words represent the same group element. The Higman embedding theorem characterises recursively presented groups.

\term{First Isomorphism Theorem}
For a group homomorphism $\phi: G \to H$, the quotient group $G/\ker(\phi)$ is isomorphic to the image of $\phi$: $G/\ker(\phi) \cong \operatorname{im}(\phi)$. The analogous statements hold for rings, where $R/\ker(\phi) \cong \operatorname{im}(\phi)$ as rings, and for modules, where $M/\ker(\phi) \cong \operatorname{im}(\phi)$ as modules. This is the first of the three isomorphism theorems, the other two following from it by applying it to suitable maps.

\term{First-order Logic}
A formal system using variables, relations, logical connectives, and quantifiers over elements of a specified domain. \par\noindent\textbf{Applications:} first-order logic underlies mathematical foundations, database queries, automated reasoning, and formal verification.

\term{Fischer--Burmeister Function}
The Fischer--Burmeister function $\phi(a,b) = \sqrt{a^2 + b^2} - a - b$ provides a smooth reformulation of the complementarity problem: $\phi(a,b) = 0$ iff $a \geq 0$, $b \geq 0$, $ab = 0$. It is used in numerical methods for variational inequalities and mathematical programs with equilibrium constraints. The function is continuously differentiable everywhere except at the origin.

\term{Fixed Point}
A fixed point of a function $f: X \to X$ is a point $x$ with $f(x) = x$. The Brouwer fixed-point theorem guarantees fixed points for continuous maps on compact convex sets in $\mathbb{R}^n$. In dynamical systems, fixed points are equilibrium solutions; their stability is determined by the eigenvalues of the Jacobian. Fixed-point iterations $x_{n+1} = f(x_n)$ converge under contraction conditions.

\term{Fixed-Point Theorem}
A fixed-point theorem guarantees a point $x$ with $f(x)=x$ under specified conditions, as in the Banach and Brouwer theorems. \par\noindent\textbf{Applications:} fixed points represent equilibria in economics, games, dynamical systems, and nonlinear equations.

\term{Flat Module}
An $R$-module $M$ is flat if tensoring with $M$ preserves exact sequences. Flatness is weaker than being free: over a local ring, a finitely presented flat module is free, but an arbitrary flat module need not be. Projective modules are flat, and flatness makes the algebraic behavior of fibers change without sudden jumps. \par\noindent\textbf{Applications:} flat modules and flat maps are central in algebraic geometry, where they describe families whose fibers vary consistently.

\term{Flatness}
Flatness means that tensoring a module with $M$ preserves exact sequences; over an integral domain it generalizes being torsion-free. \par\noindent\textbf{Applications:} flatness describes algebraic families whose fibers change without sudden jumps in algebraic geometry.

\term{Floating-Point Arithmetic}
Floating-point arithmetic represents real numbers approximately as $m\times\beta^e$ using a finite mantissa and exponent. Rounding and overflow are unavoidable and are standardized by IEEE 754. \par\noindent\textbf{Applications:} understanding floating point is essential for reliable scientific computing, graphics, finance, and machine learning.

\term{Flow}
A flow is a family of maps describing how states evolve with time; its paths are integral curves of a vector field. \par\noindent\textbf{Applications:} flows model motion, fluid transport, population change, and continuous-time control systems.

\term{Focal Point (Optics)}
A focal point of a lens or mirror is the point where parallel rays converge (or appear to diverge from). For a thin lens with focal length $f$, the lens equation is $\frac{1}{f} = \frac{1}{d_o} + \frac{1}{d_i}$ where $d_o$ and $d_i$ are object and image distances. Focal properties are described by ray transfer matrices in paraxial optics. The focal length determines the magnification and field of view.

\term{Fold Catastrophe}
A fold catastrophe is the simplest singularity in Thom's classification of elementary catastrophes. The normal form is $V(x; \lambda) = x^3 + \lambda x$, with catastrophe set given by the discriminant $27\lambda^2 + 4x^3 = 0$. At the fold, a stable and unstable branch of equilibria coalesce and annihilate. Folds appear in optics (rainbow formation), bifurcation theory, and elasticity.

\term{Foata Transform}
The Foata transform is a bijection on permutations that maps the descent set of a permutation to its cycle structure in a controlled way. Defined by Désiré Foata (1968), it preserves the Eulerian distribution and relates to the Robinson--Schensted correspondence. The transform is used in combinatorics on words, pattern avoidance, and the study of permutation statistics.

\term{Formula}
A formula is a symbolic rule relating quantities, such as the quadratic formula. \par\noindent\textbf{Applications:} formulas make models computable, communicate mathematical relationships, and support prediction and parameter estimation.

\term{Four Color Theorem}
Every planar graph has a vertex coloring using at most four colors such that adjacent vertices receive different colors. The theorem was proved in 1976 by Kenneth Appel and Wolfgang Haken, the first major theorem to be proved with substantial computer assistance. Four colors are sometimes necessary, since the complete graph $K_4$ is planar.

\term{Four-Square Theorem}
Also called Lagrange's four-square theorem: every natural number is a sum of four integer squares, $n = a_1^2 + a_2^2 + a_3^2 + a_4^2$, proved by Lagrange in 1770. Related results about sums of squares include Jacobi's four-square theorem, which counts the number of representations of $n$, and Legendre's three-square theorem, which characterises when three squares suffice.

\term{Fourier Coefficients}
For a $2\pi$-periodic function $f \in L^1([-\pi,\pi])$, the Fourier coefficients are $\hat{f}(n) = \frac{1}{2\pi} \int_{-\pi}^{\pi} f(x) e^{-inx} \, dx$ for $n \in \mathbb{Z}$. The Fourier series is $f(x) \sim \sum_{n=-\infty}^{\infty} \hat{f}(n) e^{inx}$. For $f \in L^2$, Parseval's identity states $\|f\|_2^2 = \sum |\hat{f}(n)|^2$. The coefficients measure the frequency content of $f$.

\term{Fourier Series}
The representation of a periodic function as a sum of sinusoids, $f(x) = a_0 + \sum (a_n \cos nx + b_n \sin nx)$.

\term{Fourier Transform}
The Fourier transform of $f \in L^1(\mathbb{R})$ is $\hat{f}(\xi) = \int_{-\infty}^{\infty} f(x) e^{-2\pi i x \xi} \, dx$. It extends to an isometry on $L^2(\mathbb{R})$ (Plancherel theorem). Key properties: $\widehat{f'}(\xi) = 2\pi i \xi \hat{f}(\xi)$, $\hat{\hat{f}}(x) = f(-x)$, and $\widehat{f*g} = \hat{f} \cdot \hat{g}$. The Fourier transform diagonalises convolution and is fundamental to analysis, signal processing, and quantum mechanics.

\term{Fourier--Motzkin Elimination}
Fourier--Motzkin elimination is a method for eliminating variables from a system of linear inequalities. Given inequalities in variables $x_1, \ldots, x_n$, one can eliminate $x_n$ by combining pairs of inequalities of the form $a x_n \leq b$ and $c x_n \geq d$ to produce an inequality not involving $x_n$. The procedure is the linear programming analogue of Gaussian elimination and is used in constraint projection.

\term{Fourier's Law}
Fourier's law of heat conduction states that the heat flux $\mathbf{q}$ is proportional to the negative temperature gradient: $\mathbf{q} = -k \nabla T$ where $k > 0$ is the thermal conductivity. Combined with conservation of energy, this yields the heat equation $\frac{\partial T}{\partial t} = \alpha \nabla^2 T$ where $\alpha = k/(\rho c_p)$. The law is empirical but derivable from kinetic theory.

\term{Fractal}
A fractal is a set with detail or patterns that repeat across scales, often with a non-integer dimension. Examples include the Cantor and Mandelbrot sets. \par\noindent\textbf{Applications:} fractals model coastlines, branching networks, porous materials, textures, and financial or biological patterns.

\term{Fractal Dimension}
The fractal (Hausdorff) dimension of a set $E \subset \mathbb{R}^n$ is $s = \inf\{t : \mathcal{H}^t(E) = 0\} = \sup\{t : \mathcal{H}^t(E) = \infty\}$ where $\mathcal{H}^t$ is the $t$-dimensional Hausdorff measure. For the Cantor set, $s = \log 2 / \log 3$. The Sierpiński triangle has dimension $\log 3 / \log 2$. Fractal dimensions quantify the "roughness" or "space-filling" property of self-similar sets.

\term{Frame}
A frame is a possibly redundant family of vectors that stably represents every vector in a Hilbert space. Unlike a basis, it may contain more vectors than necessary. \par\noindent\textbf{Applications:} frames enable robust signal reconstruction, image compression, sampling, and noise-resistant measurements.

\term{Frattini Subgroup}
The Frattini subgroup $\Phi(G)$ of a group $G$ is the intersection of all maximal subgroups of $G$. Equivalently, it is the set of nongenerator elements: $g \in \Phi(G)$ iff whenever $G = \langle g, S \rangle$ for some $S$, then $G = \langle S \rangle$. For a finite $p$-group, $\Phi(G) = G^p [G,G]$ (the product of $p$-th powers and commutators). The Frattini subgroup is nilpotent and characteristic.

\term{Fréchet Derivative}
The derivative of a map between normed spaces as a bounded linear operator, generalising the total derivative.

\term{Fredholm Alternative}
For a compact operator $T$ on a normed space, either the equation $(I - T)x = y$ is uniquely solvable for every $y$, or the homogeneous equations $(I - T)x = 0$ and $(I - T^*)y = 0$ have the same positive finite dimension. This dichotomy is the core of Fredholm theory and is the basis for solving integral equations with compact kernels. It reduces solvability to orthogonality conditions against solutions of the adjoint homogeneous equation.

\term{Fredholm Theory}
Fredholm theory studies equations of the form $(I-K)x=y$, where $K$ is often a compact operator. It gives a solvability alternative: an equation is solvable precisely when $y$ satisfies compatibility conditions coming from the corresponding adjoint equation. \par\noindent\textbf{Applications:} the theory helps analyze integral equations, boundary-value problems, and linearized partial differential equations.

\term{Free Group}
A free group is generated by symbols with no relations except cancellation of a symbol with its inverse. Every group is a quotient of a free group. \par\noindent\textbf{Applications:} free groups describe paths, symmetries, automata, and the presentations used in computational group theory.

\term{Free Module}
A free module has a basis, so every element has a unique finite linear combination of basis elements. Equivalently, it is a direct sum of copies of its ring. \par\noindent\textbf{Applications:} free modules provide coordinates for algebraic systems and are used in linear algebra, coding, and algebraic geometry.

\term{Frenet--Serret Formulas}
For a smooth curve $\alpha(s)$ parametrised by arc length in $\mathbb{R}^3$, the Frenet--Serret frame consists of the tangent $\mathbf{T}$, normal $\mathbf{N}$, and binormal $\mathbf{B} = \mathbf{T} \times \mathbf{N}$. The formulas are: $\mathbf{T}' = \kappa \mathbf{N}$, $\mathbf{N}' = -\kappa \mathbf{T} + \tau \mathbf{B}$, $\mathbf{B}' = -\tau \mathbf{N}$ where $\kappa$ is curvature and $\tau$ is torsion. These equations characterise curves up to rigid motion.

\term{Fremlin Measure}
The Fremlin measure is a construction in measure theory that produces a localisable measure from a premeasure on an algebra. Defined by David Fremlin, it is used in the theory of Riesz spaces and vector lattices. The Fremlin construction is essential for extending measures to larger $\sigma$-algebras while preserving locality and completeness properties.

\term{Frequency}
Frequency is the number of cycles per unit time of a periodic phenomenon; it is the reciprocal of the period. \par\noindent\textbf{Applications:} frequency is fundamental in sound, radio, vibration, image analysis, and Fourier-based signal processing.

\term{Freudenthal's Theorem (Homotopy)}
Freudenthal's theorem on the stability of homotopy groups states that for a connected CW-complex $X$, the suspension homomorphism $\pi_k(X) \to \pi_{k+1}(\Sigma X)$ is an isomorphism for $k < 2n-1$ and surjective for $k = 2n-1$ when $X$ is $(n-1)$-connected. This is the first step in the Freudenthal suspension theorem, foundational for stable homotopy theory.

\term{Frobenius Coin Problem}
The Frobenius coin problem asks for the largest integer that cannot be expressed as a non-negative integer combination of given coprime positive integers $a_1, \ldots, a_n$. For two numbers $a, b$, the Frobenius number is $g(a,b) = ab - a - b$ (Sylvester, 1884). For three or more numbers, no closed formula is known; the problem is NP-hard. The number of non-representable integers is $\frac{1}{(n-1)!} \prod a_i$ for $n=2$.

\term{Frobenius Endomorphism}
In characteristic $p > 0$, the Frobenius endomorphism is the map $x \mapsto x^p$ on a commutative ring or field. It is a ring homomorphism: $(x+y)^p = x^p + y^p$ in characteristic $p$. For a perfect field $\mathbb{F}_q$, the Frobenius generates the Galois group $\operatorname{Gal}(\mathbb{F}_q/\mathbb{F}_p) \cong \mathbb{Z}/n\mathbb{Z}$. In algebraic geometry, the Frobenius acts on cohomology and is central to the Weil conjectures.

\term{Frobenius Reciprocity}
Frobenius reciprocity relates induction and restriction of representations. If $H \leq G$, $\chi$ is a character of $H$, and $\psi$ is a character of $G$, then $\langle \operatorname{Ind}_H^G \chi, \psi \rangle_G = \langle \chi, \operatorname{Res}_H^G \psi \rangle_H$. This adjunction between induction and restriction is fundamental in representation theory and appears in many categories (modules, sheaves, Hilbert spaces).

\term{Frobenius Theorem (Distributions)}
The Frobenius theorem characterises integrable distributions on manifolds. A smooth distribution $D \subset TM$ is integrable (tangent to a foliation) iff it is involutive: $[X, Y] \in D$ for all $X, Y \in D$. In Lie group theory, a Lie subalgebra integrates to a Lie subgroup. The theorem provides the bridge between infinitesimal (Lie algebra) and global (manifold) structures.

\term{Fubini Property}
(Measure theory) The interchange of order of integration in iterated integrals for positive measurable functions, i.e. $\int\int f \, dx\, dy = \int\int f \, dy\, dx$.

\term{Fubini's Theorem}
If $f$ is integrable on a product space $X \times Y$, then $\int_{X \times Y} f \, d(\mu \times \nu) = \int_X \left(\int_Y f(x,y) \, d\nu\right) d\mu = \int_Y \left(\int_X f(x,y) \, d\mu\right) d\nu$. This justifies iterated integration for Lebesgue integrals and is foundational in analysis. Tonelli's theorem applies to non-negative functions without integrability assumptions. Fubini extends to countable products and product measures.

\term{Fuchsian Group}
A Fuchsian group is a discrete subgroup of $\operatorname{PSL}(2, \mathbb{R})$, acting on the upper half-plane $\mathbb{H}$ by Möbius transformations. The quotient $\mathbb{H}/\Gamma$ is a Riemann surface (possibly with orbifold points). Fuchsian groups generalise modular groups; the modular group $\operatorname{PSL}(2, \mathbb{Z})$ is the prototypical example. Uniformization theorem states every hyperbolic Riemann surface is a quotient by a Fuchsian group.

\term{Fuchsian Differential Equation}
A Fuchsian differential equation is a linear ODE whose singular points are all regular (solutions have at most polynomial growth near each singularity). The hypergeometric equation $z(1-z)y'' + [c - (a+b+1)z]y' - aby = 0$ has three regular singular points at $0, 1, \infty$. Riemann's $P$-function formalism classifies Fuchsian equations by their singularities and exponents.

\term{Fullerene}
A fullerene is a molecular graph representing a carbon molecule with 20 or more vertices, where each vertex has degree 3 and every face is a pentagon or hexagon. Euler's formula forces exactly 12 pentagonal faces; the number of hexagons can vary. The buckyball $C_{60}$ (truncated icosahedron) is the most symmetric fullerene. Fullerenes model nanotubes and zero-dimensional carbon allotropes.

\term{Function Space}
A space whose points are functions, such as $C(X)$ or $L^p$, with a topology or metric measuring the closeness of functions.

\term{Functional}
A map from a vector space to its scalar field, especially a continuous linear one such as integration.

\term{Functional Analysis}
Functional analysis studies infinite-dimensional vector spaces equipped with additional structure (norms, inner products, topologies). Key objects: Banach spaces (complete normed spaces), Hilbert spaces (complete inner product spaces), and operator algebras. Fundamental theorems include Hahn--Banach (extension of functionals), Banach--Steinhaus (uniform boundedness), Open Mapping, and Closed Graph theorems. Applications span PDEs, quantum mechanics, and harmonic analysis.

\term{Functional Calculus}
The construction of functions of an operator, such as $f(A)$, via spectral measures or power series.

\term{Functional Determinant}
The functional determinant generalises the matrix determinant to operators. For a differential operator $L$ on a domain with boundary conditions, $\det L$ is defined via zeta-function regularisation: $\det L = \exp(-\zeta_L'(0))$ where $\zeta_L(s) = \sum \lambda_n^{-s}$ sums over eigenvalues. The determinant appears in path integrals, spectral geometry, and the study of quantum anomalies.

\term{Functor}
A structure-preserving map between categories, sending objects to objects and morphisms to morphisms while preserving composition and identities.

\term{Fundamental Group}
The fundamental group $\pi_1(X, x_0)$ consists of homotopy classes of loops based at $x_0$, with composition given by concatenation. For the circle, $\pi_1(S^1) \cong \mathbb{Z}$ (winding number). The fundamental group detects topological features: it is trivial for simply connected spaces and non-trivial for spaces with "holes." Covering space theory establishes a correspondence between subgroups of $\pi_1$ and covering spaces.

\term{Fundamental Solution}
A distribution solving a linear PDE with a delta-function source, from which general solutions are built by convolution.

\term{Fundamental Theorem of Algebra}
Every non-constant polynomial with complex coefficients has at least one complex root. Equivalently, $\mathbb{C}$ is algebraically closed: every polynomial of degree $n$ has exactly $n$ roots counting multiplicity. Proofs exist via complex analysis (Liouville's theorem), algebraic topology (degree theory), and real analysis (infimum argument). The theorem was first stated by Gauss (1799) and is foundational to algebra and analysis.

\term{Fundamental Theorem of Algebraic Topology}
Homology (and cohomology) is a functor from the homotopy category of topological spaces to the category of groups: a continuous map $f: X \to Y$ induces homomorphisms $f_*: H_n(X) \to H_n(Y)$, and a homotopy equivalence induces an isomorphism for every $n$. Spaces of the same homotopy type therefore have isomorphic homology groups, making homology a topological invariant of homotopy type.

\term{Fundamental Theorem of Arithmetic}
Every integer greater than 1 factors uniquely, up to the order of the factors, into a product of primes. This uniqueness is the reason primes are described as the atoms of integer multiplication. The theorem also underlies the Euclidean algorithm and the computation of greatest common divisors.

\term{Fundamental Theorem of Calculus}
If $f$ is continuous on $[a,b]$, then $F(x) = \int_a^x f(t) \, dt$ is differentiable with $F'(x) = f(x)$ (Part 1). If $F$ is any antiderivative of $f$, then $\int_a^b f(x) \, dx = F(b) - F(a)$ (Part 2). This theorem links differentiation and integration, the two central operations of calculus. It generalises to manifolds via Stokes' theorem: $\int_M d\omega = \int_{\partial M} \omega$.

\term{Fundamental Theorem of Finite Abelian Groups}
Every finite abelian group is isomorphic to a direct product of cyclic groups of prime-power order, $G \cong C_{p_1^{a_1}} \times \cdots \times C_{p_k^{a_k}}$, and the decomposition is unique up to the order of the factors. Equivalently, a finite abelian group is the direct product of its Sylow $p$-subgroups, each of which splits into cyclic factors; this is the primary decomposition. The theorem classifies all finite abelian groups.

\term{Fundamental Theorem of Galois Theory}
For a finite Galois extension $L/K$ with group $G = \operatorname{Gal}(L/K)$, there is an inclusion-reversing bijection between intermediate fields $K \subseteq M \subseteq L$ and subgroups $H \leq G$: $M \mapsto \operatorname{Gal}(L/M)$ and $H \mapsto L^H$ (fixed field). Normal subgroups correspond to Galois extensions. This theorem is the cornerstone of Galois theory, linking field theory to group theory.

\term{Fundamental Theorem of Lie Algebras}
A finite-dimensional Lie algebra over $\mathbb{C}$ decomposes as $\mathfrak{g} = \mathfrak{s} \ltimes \mathfrak{r}$ where $\mathfrak{s}$ is semisimple (direct sum of simple algebras) and $\mathfrak{r}$ is the solvable radical (maximal solvable ideal). Cartan's criterion characterises semisimplicity by the non-degeneracy of the Killing form $B(x,y) = \operatorname{tr}(\operatorname{ad}_x \operatorname{ad}_y)$. This structural theorem parallels the Jordan--Hölder theorem for groups.

\term{Fundamental Theorem of Linear Algebra}
For an $m \times n$ matrix $A$, the four fundamental subspaces satisfy: (1) $\dim \operatorname{Nul}(A) + \dim \operatorname{Col}(A) = n$ (rank-nullity); (2) $\operatorname{Nul}(A) = \operatorname{Col}(A^T)^\perp$; (3) $\operatorname{Nul}(A^T) = \operatorname{Col}(A)^\perp$; (4) $\operatorname{rank}(A) = \operatorname{rank}(A^T)$. This theorem, due to Gilbert Strang, provides the geometric foundation for matrix analysis and is essential for understanding systems of linear equations.

\term{Fundamental Theorem of Symmetric Functions}
Every symmetric polynomial in $x_1, \ldots, x_n$ can be uniquely expressed as a polynomial in the elementary symmetric polynomials $e_k = \sum_{i_1 < \cdots < i_k} x_{i_1} \cdots x_{i_k}$. The ring of symmetric polynomials is $\mathbb{Z}[e_1, \ldots, e_n]$. This theorem underlies the theory of resultants, discriminants, and invariant theory. It is also the algebraic basis for expressing coefficients of polynomials in terms of their roots.

\term{Future Cone}
The future light cone at a point $p$ in Minkowski spacetime is the set $\{q : q - p \text{ is future-directed timelike or null}\}$. In coordinates with signature $(-,+,+,+)$, the future cone is $\{t > 0, t^2 > x^2 + y^2 + z^2\}$. The future cone defines the causal structure of spacetime: events in the future cone can be influenced by $p$. Light cones generalise to curved spacetimes in general relativity.
\term{Coequalizer}
The dual notion to the equalizer. The coequalizer of $f, g: A \to B$ is an object $C$ with $q: B \to C$ such that $q \circ f = q \circ g$, universal among such morphisms. In sets, the coequalizer is the quotient of $B$ by the equivalence relation generated by $f(a) \sim g(a)$ for all $a \in A$.
\term{Commutative Diagram}
A diagram in a category in which any two paths with the same endpoints give the same composed morphism. Commutative diagrams are the primary language for expressing universal properties, naturality, and exactness in algebra and topology.
\term{Feasible Region}
In optimization, the set of all points satisfying the constraints of a problem. For a linear program, the feasible region is a convex polyhedron, and the optimum occurs at a vertex. In nonlinear programming, the feasible region may be a more general subset of $\mathbb{R}^n$.
\term{Finite Difference}
An approximation to a derivative using values of a function at discrete points: the forward difference $\Delta f(x) = f(x+h) - f(x)$, the backward difference, and the central difference. Finite difference methods replace differential equations by recurrence relations on grids. They are foundational to numerical analysis and the discretization of PDEs.
